% ============================================================
%  Appendix I  Joint TC-Chunking x TC-Metadata indexing
%  % ============================================================
%  Appendix I  Joint TC-Chunking x TC-Metadata indexing
%  % ============================================================
%  Appendix I  Joint TC-Chunking x TC-Metadata indexing
%  % ============================================================
%  Appendix I  Joint TC-Chunking x TC-Metadata indexing
%  \input{I_joint_analysis} in main paper
% ============================================================

\section{Joint TC-Chunking $\times$ TC-Metadata Indexing}
\label{sec:appendix-joint}

For completeness we report the joint optimisation of the two index-side
axes (chunking and metadata) that was excluded from the main results
because it does not strictly improve on the single-axis TC-Metadata
variant under any retriever. We describe the setup, the held-out
numbers, and the structural reasons for the non-significance.

\paragraph{Methodology.} The strategy set is the Cartesian product
$\mathcal{S} = \mathcal{S}_{\text{chunk}} \times \mathcal{S}_{\text{meta}}$
of size $14 \times 14 = 196$. For each pair $(c, m) \in \mathcal{S}$ the
\textit{Les Débats} index is the $c$-chunked collection re-queried
under the metadata predicate or reranker $m$; newspaper retrieval is
left at the global default. Inference follows the shared protocol of
Section~\ref{sec:indexing-methodology}, with the only change that each
cluster stores a \emph{pair} $s^{*}(c) = (c^{*}, m^{*})$ rather than a
single index choice.
Table~\ref{tab:joint-wins-summary}.

\paragraph{Held-out results.}
Table~\ref{tab:tc-joint-vs-metadata} reports the joint variant
side-by-side with TC-Metadata on the held-out test set
($n_{\text{test}}{=}438$, 5 seeds). Under the Naive Retriever the joint
pair adds $+8.8$~\% points Coverage@3 and $+8.6$~\% points Recall@3
($p<0.05$), marginally below TC-Metadata ($+8.9$ / $+9.0$~\% points).
Under UMS the gains are $+6.0$~\% points Cov@3 / $+6.0$~\% points Rec@3
($p<0.05$), matching TC-Metadata to within a tenth of a point. QRe and
QRe\,\&\,UMS show no significant improvement under either indexing
variant (ns). The joint variant therefore never strictly dominates
TC-Metadata.

\begin{table}[htbp]
\centering
\small
\setlength{\tabcolsep}{4pt}
\renewcommand{\arraystretch}{1.1}
\resizebox{\columnwidth}{!}{%
\begin{tabular}{ll cc cc}
\toprule
\textbf{Retriever} & \textbf{Indexing strategy}
  & \multicolumn{2}{c}{\textbf{Coverage (\%)}}
  & \multicolumn{2}{c}{\textbf{Recall (\%)}} \\
\cmidrule(lr){3-4} \cmidrule(lr){5-6}
 & & @3 & @5 & @3 & @5 \\
\midrule
\multirow{3}{*}{Naive}
  & Global Baseline      & 36.9 & 45.4 & 36.9 & 45.6 \\
  & \;+\,TC-Metadata     & 45.8$^{*}$  & 55.2$^{*}$  & 45.9$^{*}$   & 55.3$^{*}$ \\
  & \;+\,TC-Joint        & 45.7$^{*}$  & 54.7$^{*}$  & 45.5$^{**}$  & 54.5$^{*}$ \\
\midrule
\multirow{3}{*}{QRe}
  & Global Baseline      & 48.5 & 58.2 & 48.7 & 58.4 \\
  & \;+\,TC-Metadata     & 48.5$^{\text{ns}}$ & 58.2$^{\text{ns}}$ & 48.7$^{\text{ns}}$ & 58.4$^{\text{ns}}$ \\
  & \;+\,TC-Joint        & 48.3$^{\text{ns}}$ & 57.6$^{\text{ns}}$ & 48.2$^{\text{ns}}$ & 57.4$^{\text{ns}}$ \\
\midrule
\multirow{3}{*}{UMS}
  & Global Baseline      & 46.8 & 58.7 & 47.0 & 58.8 \\
  & \;+\,TC-Metadata     & 52.8$^{*}$ & 63.2$^{*}$ & 53.0$^{*}$ & 63.4$^{*}$ \\
  & \;+\,TC-Joint        & 52.8$^{*}$ & 63.1$^{*}$ & 53.0$^{*}$ & 63.4$^{*}$ \\
\midrule
\multirow{3}{*}{QRe \& UMS}
  & Global Baseline      & 54.7 & 64.5 & 54.9 & 64.7 \\
  & \;+\,TC-Metadata     & 54.7$^{\text{ns}}$ & 64.5$^{\text{ns}}$ & 54.9$^{\text{ns}}$ & 64.7$^{\text{ns}}$ \\
  & \;+\,TC-Joint        & 54.7$^{\text{ns}}$ & 64.5$^{\text{ns}}$ & 54.9$^{\text{ns}}$ & 64.7$^{\text{ns}}$ \\
\bottomrule
\end{tabular}
}
\caption{Joint chunking$\times$metadata variant compared to the
single-axis TC-Metadata variant on the held-out test set
($n_{\text{test}}{=}438$, 5 seeds). Each test question is routed to
its nearest question cluster, and the chunking and metadata strategies
used to retrieve its chunks are those found optimal for that cluster
on the training split. Under \textit{+\,TC-Metadata}, the cluster
picks its best metadata configuration (out of 14, chunking fixed at
\texttt{10000\_hier.}); under \textit{+\,TC-Joint}, the cluster picks
its best (chunking, metadata) pair jointly from the full
$14{\times}14{=}196$-pair grid. Both variants are therefore
per-cluster, i.e.\ effectively per-question.
$^{*}p<0.05$, $^{**}p<0.01$ (Bonferroni paired $t$-test);
ns = not significant.}
\label{tab:tc-joint-vs-metadata}
\end{table}

\paragraph{Why the joint variant does not improve on TC-Metadata.}
Two effects combine to flatten the joint gain.
\textbf{(i) Winner's curse over a 196-pair grid.} The cluster-level
selector picks the best pair on the training split out of 196
candidates rather than 14, so the selection variance grows roughly with
$\log(|\mathcal{S}|)$ at fixed training-cluster size. With $\sim$20
training clusters and 5 test seeds, the inflated train-time pair score
does not transfer fully to held-out questions, eroding any
configuration-level advantage the joint search might have had.
\textbf{(ii) Correlated gain channels.} Chunking and metadata strategies
on \textit{Les Débats} both operate on the same failure mode:
suppressing the parliamentary noise that crowds out the top-$k$
(short \texttt{fragment} / \texttt{vote\_list} chunks for chunking;
\texttt{exclude\_noisy} predicates and type-boost rerankers for
metadata). Once either axis has filtered or down-weighted those chunks,
the marginal contribution of the other axis on the same questions is
near zero, so the joint pair recovers the TC-Metadata gain rather than
stacking on top of it. The two effects together are consistent with the
empirical pattern, joint matches TC-Metadata where TC-Metadata helps
(Naive, UMS), and stays at noise level where TC-Metadata is already
neutral (QRe, QRe\,\&\,UMS).
 in main paper
% ============================================================

\section{Joint TC-Chunking $\times$ TC-Metadata Indexing}
\label{sec:appendix-joint}

For completeness we report the joint optimisation of the two index-side
axes (chunking and metadata) that was excluded from the main results
because it does not strictly improve on the single-axis TC-Metadata
variant under any retriever. We describe the setup, the held-out
numbers, and the structural reasons for the non-significance.

\paragraph{Methodology.} The strategy set is the Cartesian product
$\mathcal{S} = \mathcal{S}_{\text{chunk}} \times \mathcal{S}_{\text{meta}}$
of size $14 \times 14 = 196$. For each pair $(c, m) \in \mathcal{S}$ the
\textit{Les Débats} index is the $c$-chunked collection re-queried
under the metadata predicate or reranker $m$; newspaper retrieval is
left at the global default. Inference follows the shared protocol of
Section~\ref{sec:indexing-methodology}, with the only change that each
cluster stores a \emph{pair} $s^{*}(c) = (c^{*}, m^{*})$ rather than a
single index choice.
Table~\ref{tab:joint-wins-summary}.

\paragraph{Held-out results.}
Table~\ref{tab:tc-joint-vs-metadata} reports the joint variant
side-by-side with TC-Metadata on the held-out test set
($n_{\text{test}}{=}438$, 5 seeds). Under the Naive Retriever the joint
pair adds $+8.8$~\% points Coverage@3 and $+8.6$~\% points Recall@3
($p<0.05$), marginally below TC-Metadata ($+8.9$ / $+9.0$~\% points).
Under UMS the gains are $+6.0$~\% points Cov@3 / $+6.0$~\% points Rec@3
($p<0.05$), matching TC-Metadata to within a tenth of a point. QRe and
QRe\,\&\,UMS show no significant improvement under either indexing
variant (ns). The joint variant therefore never strictly dominates
TC-Metadata.

\begin{table}[htbp]
\centering
\small
\setlength{\tabcolsep}{4pt}
\renewcommand{\arraystretch}{1.1}
\resizebox{\columnwidth}{!}{%
\begin{tabular}{ll cc cc}
\toprule
\textbf{Retriever} & \textbf{Indexing strategy}
  & \multicolumn{2}{c}{\textbf{Coverage (\%)}}
  & \multicolumn{2}{c}{\textbf{Recall (\%)}} \\
\cmidrule(lr){3-4} \cmidrule(lr){5-6}
 & & @3 & @5 & @3 & @5 \\
\midrule
\multirow{3}{*}{Naive}
  & Global Baseline      & 36.9 & 45.4 & 36.9 & 45.6 \\
  & \;+\,TC-Metadata     & 45.8$^{*}$  & 55.2$^{*}$  & 45.9$^{*}$   & 55.3$^{*}$ \\
  & \;+\,TC-Joint        & 45.7$^{*}$  & 54.7$^{*}$  & 45.5$^{**}$  & 54.5$^{*}$ \\
\midrule
\multirow{3}{*}{QRe}
  & Global Baseline      & 48.5 & 58.2 & 48.7 & 58.4 \\
  & \;+\,TC-Metadata     & 48.5$^{\text{ns}}$ & 58.2$^{\text{ns}}$ & 48.7$^{\text{ns}}$ & 58.4$^{\text{ns}}$ \\
  & \;+\,TC-Joint        & 48.3$^{\text{ns}}$ & 57.6$^{\text{ns}}$ & 48.2$^{\text{ns}}$ & 57.4$^{\text{ns}}$ \\
\midrule
\multirow{3}{*}{UMS}
  & Global Baseline      & 46.8 & 58.7 & 47.0 & 58.8 \\
  & \;+\,TC-Metadata     & 52.8$^{*}$ & 63.2$^{*}$ & 53.0$^{*}$ & 63.4$^{*}$ \\
  & \;+\,TC-Joint        & 52.8$^{*}$ & 63.1$^{*}$ & 53.0$^{*}$ & 63.4$^{*}$ \\
\midrule
\multirow{3}{*}{QRe \& UMS}
  & Global Baseline      & 54.7 & 64.5 & 54.9 & 64.7 \\
  & \;+\,TC-Metadata     & 54.7$^{\text{ns}}$ & 64.5$^{\text{ns}}$ & 54.9$^{\text{ns}}$ & 64.7$^{\text{ns}}$ \\
  & \;+\,TC-Joint        & 54.7$^{\text{ns}}$ & 64.5$^{\text{ns}}$ & 54.9$^{\text{ns}}$ & 64.7$^{\text{ns}}$ \\
\bottomrule
\end{tabular}
}
\caption{Joint chunking$\times$metadata variant compared to the
single-axis TC-Metadata variant on the held-out test set
($n_{\text{test}}{=}438$, 5 seeds). Each test question is routed to
its nearest question cluster, and the chunking and metadata strategies
used to retrieve its chunks are those found optimal for that cluster
on the training split. Under \textit{+\,TC-Metadata}, the cluster
picks its best metadata configuration (out of 14, chunking fixed at
\texttt{10000\_hier.}); under \textit{+\,TC-Joint}, the cluster picks
its best (chunking, metadata) pair jointly from the full
$14{\times}14{=}196$-pair grid. Both variants are therefore
per-cluster, i.e.\ effectively per-question.
$^{*}p<0.05$, $^{**}p<0.01$ (Bonferroni paired $t$-test);
ns = not significant.}
\label{tab:tc-joint-vs-metadata}
\end{table}

\paragraph{Why the joint variant does not improve on TC-Metadata.}
Two effects combine to flatten the joint gain.
\textbf{(i) Winner's curse over a 196-pair grid.} The cluster-level
selector picks the best pair on the training split out of 196
candidates rather than 14, so the selection variance grows roughly with
$\log(|\mathcal{S}|)$ at fixed training-cluster size. With $\sim$20
training clusters and 5 test seeds, the inflated train-time pair score
does not transfer fully to held-out questions, eroding any
configuration-level advantage the joint search might have had.
\textbf{(ii) Correlated gain channels.} Chunking and metadata strategies
on \textit{Les Débats} both operate on the same failure mode:
suppressing the parliamentary noise that crowds out the top-$k$
(short \texttt{fragment} / \texttt{vote\_list} chunks for chunking;
\texttt{exclude\_noisy} predicates and type-boost rerankers for
metadata). Once either axis has filtered or down-weighted those chunks,
the marginal contribution of the other axis on the same questions is
near zero, so the joint pair recovers the TC-Metadata gain rather than
stacking on top of it. The two effects together are consistent with the
empirical pattern, joint matches TC-Metadata where TC-Metadata helps
(Naive, UMS), and stays at noise level where TC-Metadata is already
neutral (QRe, QRe\,\&\,UMS).
 in main paper
% ============================================================

\section{Joint TC-Chunking $\times$ TC-Metadata Indexing}
\label{sec:appendix-joint}

For completeness we report the joint optimisation of the two index-side
axes (chunking and metadata) that was excluded from the main results
because it does not strictly improve on the single-axis TC-Metadata
variant under any retriever. We describe the setup, the held-out
numbers, and the structural reasons for the non-significance.

\paragraph{Methodology.} The strategy set is the Cartesian product
$\mathcal{S} = \mathcal{S}_{\text{chunk}} \times \mathcal{S}_{\text{meta}}$
of size $14 \times 14 = 196$. For each pair $(c, m) \in \mathcal{S}$ the
\textit{Les Débats} index is the $c$-chunked collection re-queried
under the metadata predicate or reranker $m$; newspaper retrieval is
left at the global default. Inference follows the shared protocol of
Section~\ref{sec:indexing-methodology}, with the only change that each
cluster stores a \emph{pair} $s^{*}(c) = (c^{*}, m^{*})$ rather than a
single index choice.
Table~\ref{tab:joint-wins-summary}.

\paragraph{Held-out results.}
Table~\ref{tab:tc-joint-vs-metadata} reports the joint variant
side-by-side with TC-Metadata on the held-out test set
($n_{\text{test}}{=}438$, 5 seeds). Under the Naive Retriever the joint
pair adds $+8.8$~\% points Coverage@3 and $+8.6$~\% points Recall@3
($p<0.05$), marginally below TC-Metadata ($+8.9$ / $+9.0$~\% points).
Under UMS the gains are $+6.0$~\% points Cov@3 / $+6.0$~\% points Rec@3
($p<0.05$), matching TC-Metadata to within a tenth of a point. QRe and
QRe\,\&\,UMS show no significant improvement under either indexing
variant (ns). The joint variant therefore never strictly dominates
TC-Metadata.

\begin{table}[htbp]
\centering
\small
\setlength{\tabcolsep}{4pt}
\renewcommand{\arraystretch}{1.1}
\resizebox{\columnwidth}{!}{%
\begin{tabular}{ll cc cc}
\toprule
\textbf{Retriever} & \textbf{Indexing strategy}
  & \multicolumn{2}{c}{\textbf{Coverage (\%)}}
  & \multicolumn{2}{c}{\textbf{Recall (\%)}} \\
\cmidrule(lr){3-4} \cmidrule(lr){5-6}
 & & @3 & @5 & @3 & @5 \\
\midrule
\multirow{3}{*}{Naive}
  & Global Baseline      & 36.9 & 45.4 & 36.9 & 45.6 \\
  & \;+\,TC-Metadata     & 45.8$^{*}$  & 55.2$^{*}$  & 45.9$^{*}$   & 55.3$^{*}$ \\
  & \;+\,TC-Joint        & 45.7$^{*}$  & 54.7$^{*}$  & 45.5$^{**}$  & 54.5$^{*}$ \\
\midrule
\multirow{3}{*}{QRe}
  & Global Baseline      & 48.5 & 58.2 & 48.7 & 58.4 \\
  & \;+\,TC-Metadata     & 48.5$^{\text{ns}}$ & 58.2$^{\text{ns}}$ & 48.7$^{\text{ns}}$ & 58.4$^{\text{ns}}$ \\
  & \;+\,TC-Joint        & 48.3$^{\text{ns}}$ & 57.6$^{\text{ns}}$ & 48.2$^{\text{ns}}$ & 57.4$^{\text{ns}}$ \\
\midrule
\multirow{3}{*}{UMS}
  & Global Baseline      & 46.8 & 58.7 & 47.0 & 58.8 \\
  & \;+\,TC-Metadata     & 52.8$^{*}$ & 63.2$^{*}$ & 53.0$^{*}$ & 63.4$^{*}$ \\
  & \;+\,TC-Joint        & 52.8$^{*}$ & 63.1$^{*}$ & 53.0$^{*}$ & 63.4$^{*}$ \\
\midrule
\multirow{3}{*}{QRe \& UMS}
  & Global Baseline      & 54.7 & 64.5 & 54.9 & 64.7 \\
  & \;+\,TC-Metadata     & 54.7$^{\text{ns}}$ & 64.5$^{\text{ns}}$ & 54.9$^{\text{ns}}$ & 64.7$^{\text{ns}}$ \\
  & \;+\,TC-Joint        & 54.7$^{\text{ns}}$ & 64.5$^{\text{ns}}$ & 54.9$^{\text{ns}}$ & 64.7$^{\text{ns}}$ \\
\bottomrule
\end{tabular}
}
\caption{Joint chunking$\times$metadata variant compared to the
single-axis TC-Metadata variant on the held-out test set
($n_{\text{test}}{=}438$, 5 seeds). Each test question is routed to
its nearest question cluster, and the chunking and metadata strategies
used to retrieve its chunks are those found optimal for that cluster
on the training split. Under \textit{+\,TC-Metadata}, the cluster
picks its best metadata configuration (out of 14, chunking fixed at
\texttt{10000\_hier.}); under \textit{+\,TC-Joint}, the cluster picks
its best (chunking, metadata) pair jointly from the full
$14{\times}14{=}196$-pair grid. Both variants are therefore
per-cluster, i.e.\ effectively per-question.
$^{*}p<0.05$, $^{**}p<0.01$ (Bonferroni paired $t$-test);
ns = not significant.}
\label{tab:tc-joint-vs-metadata}
\end{table}

\paragraph{Why the joint variant does not improve on TC-Metadata.}
Two effects combine to flatten the joint gain.
\textbf{(i) Winner's curse over a 196-pair grid.} The cluster-level
selector picks the best pair on the training split out of 196
candidates rather than 14, so the selection variance grows roughly with
$\log(|\mathcal{S}|)$ at fixed training-cluster size. With $\sim$20
training clusters and 5 test seeds, the inflated train-time pair score
does not transfer fully to held-out questions, eroding any
configuration-level advantage the joint search might have had.
\textbf{(ii) Correlated gain channels.} Chunking and metadata strategies
on \textit{Les Débats} both operate on the same failure mode:
suppressing the parliamentary noise that crowds out the top-$k$
(short \texttt{fragment} / \texttt{vote\_list} chunks for chunking;
\texttt{exclude\_noisy} predicates and type-boost rerankers for
metadata). Once either axis has filtered or down-weighted those chunks,
the marginal contribution of the other axis on the same questions is
near zero, so the joint pair recovers the TC-Metadata gain rather than
stacking on top of it. The two effects together are consistent with the
empirical pattern, joint matches TC-Metadata where TC-Metadata helps
(Naive, UMS), and stays at noise level where TC-Metadata is already
neutral (QRe, QRe\,\&\,UMS).
 in main paper
% ============================================================

\section{Joint TC-Chunking $\times$ TC-Metadata Indexing}
\label{sec:appendix-joint}

For completeness we report the joint optimisation of the two index-side
axes (chunking and metadata) that was excluded from the main results
because it does not strictly improve on the single-axis TC-Metadata
variant under any retriever. We describe the setup, the held-out
numbers, and the structural reasons for the non-significance.

\paragraph{Methodology.} The strategy set is the Cartesian product
$\mathcal{S} = \mathcal{S}_{\text{chunk}} \times \mathcal{S}_{\text{meta}}$
of size $14 \times 14 = 196$. For each pair $(c, m) \in \mathcal{S}$ the
\textit{Les Débats} index is the $c$-chunked collection re-queried
under the metadata predicate or reranker $m$; newspaper retrieval is
left at the global default. Inference follows the shared protocol of
Section~\ref{sec:indexing-methodology}, with the only change that each
cluster stores a \emph{pair} $s^{*}(c) = (c^{*}, m^{*})$ rather than a
single index choice.
Table~\ref{tab:joint-wins-summary}.

\paragraph{Held-out results.}
Table~\ref{tab:tc-joint-vs-metadata} reports the joint variant
side-by-side with TC-Metadata on the held-out test set
($n_{\text{test}}{=}438$, 5 seeds). Under the Naive Retriever the joint
pair adds $+8.8$~\% points Coverage@3 and $+8.6$~\% points Recall@3
($p<0.05$), marginally below TC-Metadata ($+8.9$ / $+9.0$~\% points).
Under UMS the gains are $+6.0$~\% points Cov@3 / $+6.0$~\% points Rec@3
($p<0.05$), matching TC-Metadata to within a tenth of a point. QRe and
QRe\,\&\,UMS show no significant improvement under either indexing
variant (ns). The joint variant therefore never strictly dominates
TC-Metadata.

\begin{table}[htbp]
\centering
\small
\setlength{\tabcolsep}{4pt}
\renewcommand{\arraystretch}{1.1}
\resizebox{\columnwidth}{!}{%
\begin{tabular}{ll cc cc}
\toprule
\textbf{Retriever} & \textbf{Indexing strategy}
  & \multicolumn{2}{c}{\textbf{Coverage (\%)}}
  & \multicolumn{2}{c}{\textbf{Recall (\%)}} \\
\cmidrule(lr){3-4} \cmidrule(lr){5-6}
 & & @3 & @5 & @3 & @5 \\
\midrule
\multirow{3}{*}{Naive}
  & Global Baseline      & 36.9 & 45.4 & 36.9 & 45.6 \\
  & \;+\,TC-Metadata     & 45.8$^{*}$  & 55.2$^{*}$  & 45.9$^{*}$   & 55.3$^{*}$ \\
  & \;+\,TC-Joint        & 45.7$^{*}$  & 54.7$^{*}$  & 45.5$^{**}$  & 54.5$^{*}$ \\
\midrule
\multirow{3}{*}{QRe}
  & Global Baseline      & 48.5 & 58.2 & 48.7 & 58.4 \\
  & \;+\,TC-Metadata     & 48.5$^{\text{ns}}$ & 58.2$^{\text{ns}}$ & 48.7$^{\text{ns}}$ & 58.4$^{\text{ns}}$ \\
  & \;+\,TC-Joint        & 48.3$^{\text{ns}}$ & 57.6$^{\text{ns}}$ & 48.2$^{\text{ns}}$ & 57.4$^{\text{ns}}$ \\
\midrule
\multirow{3}{*}{UMS}
  & Global Baseline      & 46.8 & 58.7 & 47.0 & 58.8 \\
  & \;+\,TC-Metadata     & 52.8$^{*}$ & 63.2$^{*}$ & 53.0$^{*}$ & 63.4$^{*}$ \\
  & \;+\,TC-Joint        & 52.8$^{*}$ & 63.1$^{*}$ & 53.0$^{*}$ & 63.4$^{*}$ \\
\midrule
\multirow{3}{*}{QRe \& UMS}
  & Global Baseline      & 54.7 & 64.5 & 54.9 & 64.7 \\
  & \;+\,TC-Metadata     & 54.7$^{\text{ns}}$ & 64.5$^{\text{ns}}$ & 54.9$^{\text{ns}}$ & 64.7$^{\text{ns}}$ \\
  & \;+\,TC-Joint        & 54.7$^{\text{ns}}$ & 64.5$^{\text{ns}}$ & 54.9$^{\text{ns}}$ & 64.7$^{\text{ns}}$ \\
\bottomrule
\end{tabular}
}
\caption{Joint chunking$\times$metadata variant compared to the
single-axis TC-Metadata variant on the held-out test set
($n_{\text{test}}{=}438$, 5 seeds). Each test question is routed to
its nearest question cluster, and the chunking and metadata strategies
used to retrieve its chunks are those found optimal for that cluster
on the training split. Under \textit{+\,TC-Metadata}, the cluster
picks its best metadata configuration (out of 14, chunking fixed at
\texttt{10000\_hier.}); under \textit{+\,TC-Joint}, the cluster picks
its best (chunking, metadata) pair jointly from the full
$14{\times}14{=}196$-pair grid. Both variants are therefore
per-cluster, i.e.\ effectively per-question.
$^{*}p<0.05$, $^{**}p<0.01$ (Bonferroni paired $t$-test);
ns = not significant.}
\label{tab:tc-joint-vs-metadata}
\end{table}

\paragraph{Why the joint variant does not improve on TC-Metadata.}
Two effects combine to flatten the joint gain.
\textbf{(i) Winner's curse over a 196-pair grid.} The cluster-level
selector picks the best pair on the training split out of 196
candidates rather than 14, so the selection variance grows roughly with
$\log(|\mathcal{S}|)$ at fixed training-cluster size. With $\sim$20
training clusters and 5 test seeds, the inflated train-time pair score
does not transfer fully to held-out questions, eroding any
configuration-level advantage the joint search might have had.
\textbf{(ii) Correlated gain channels.} Chunking and metadata strategies
on \textit{Les Débats} both operate on the same failure mode:
suppressing the parliamentary noise that crowds out the top-$k$
(short \texttt{fragment} / \texttt{vote\_list} chunks for chunking;
\texttt{exclude\_noisy} predicates and type-boost rerankers for
metadata). Once either axis has filtered or down-weighted those chunks,
the marginal contribution of the other axis on the same questions is
near zero, so the joint pair recovers the TC-Metadata gain rather than
stacking on top of it. The two effects together are consistent with the
empirical pattern, joint matches TC-Metadata where TC-Metadata helps
(Naive, UMS), and stays at noise level where TC-Metadata is already
neutral (QRe, QRe\,\&\,UMS).
